\section{Related Work}
\label{sec:related}

\subsection{Scene Text & Meter Recognition}
Vision-based text recognition models such as SVTR~\cite{du2022svtr} and SVTRv2~\cite{svtrv22025} utilize single-visual Transformers with local/global mixing and height-progressive merging. While highly efficient, direct CTC recognition on SVTR features remains susceptible to boundary noise and severe character confusion between geometrically similar digits (e.g. $8 \leftrightarrow 9, 3 \leftrightarrow 8, 0 \leftrightarrow 6$). Furthermore, PerturbCTC~\cite{perturbctc2026} introduces convolutional feature perturbation during training to regularize CTC alignment, yet lacks explicit spatial-coordinate grounding at inference time.

\subsection{Non-Autoregressive Sequence Refinement}
Levenshtein Transformer models and LevOCR treat sequence prediction as iterative deletion and insertion operations. However, in visual meter recognition with repetitive digit patterns (e.g. \texttt{00000.0}), unconstrained cross-attention suffers from attention scattering across identical visual glyphs. EditCTC overcomes this limitation via analytical Gaussian Spatial Priors ($\text{ARCH-4}$) and Continuous 4D Uncertainty Injection ($\text{ARCH-4C}$).

\subsection{Spec-Driven Autonomous AI & Knowledge Compilation}
The ARIA framework~\cite{chen2025spec} established the standard for Spec-Driven AI systems, decomposing complex engineering tasks into clean interoperable layers. Furthermore, WikiSkill~\cite{wikiskill2026} demonstrated that maintaining a persistent, compounding three-layer knowledge base (\texttt{raw/}, \texttt{wiki/}, \texttt{skills/}) prevents knowledge fragmentation and accelerates architectural discovery across successive model revisions. We follow these principles in organizing the research and empirical validation of EditCTC.
